\documentclass[11pt]{article}
\usepackage[margin=1in]{geometry}
\usepackage{amsmath, amssymb, amsthm}
\usepackage{hyperref}

\title{An Adaptive Funnel-Based Stopping Strategy with No-Realized-Loss Exits}
\author{Louis Charistias Saragih}
\date{May 2026}

\newtheorem{theorem}{Theorem}
\newtheorem{definition}{Definition}
\newtheorem{remark}{Remark}

\begin{document}
\maketitle

\begin{abstract}
We propose a threshold-based stopping strategy for trading under drift and volatility. The method compares a log-price path against a drift-adjusted predictive funnel and exits a position only when three conditions are simultaneously satisfied: the price crosses the upper funnel boundary, local momentum weakens, and the log-return exceeds a transaction-cost-adjusted profit threshold. This gives a rigorous no-realized-loss property for completed algorithmic exits. The guarantee does not imply that the strategy is risk-free: mark-to-market wealth may decline while holding, capital may be locked for long periods, and forced exits may realize losses. The framework is formulated using discrete-time stochastic processes, stopping times, and elementary statistical estimation.
\end{abstract}

\section{Introduction}

A common trading intuition is to hold a stock while it is below or near the purchase price, and to sell only after a sufficiently favorable move once short-term upward momentum has weakened. The purpose of this paper is to formalize this idea as a rule-based stopping strategy rather than to claim global optimality.

The strategy combines three components:
\begin{enumerate}
    \item a drift-adjusted deviation statistic measuring how far the log-price has moved relative to its expected path;
    \item a local momentum statistic detecting whether recent upward movement has stalled;
    \item an explicit profit threshold ensuring that completed algorithmic exits do not realize a loss after transaction costs.
\end{enumerate}

The central point is that the strategy can impose a \emph{no voluntary realized loss} constraint, but it cannot eliminate economic risk. While a position remains open, the mark-to-market value of the position can fall. Thus the strategy transforms realized-loss risk into holding-time risk, opportunity cost, and mark-to-market drawdown.

\section{Price Model}

Let \((S_t)_{t\geq 0}\) denote the asset price at discrete times. Define the log-price
\[
X_t=\log S_t.
\]

For the baseline model, assume
\[
X_{t+1}=X_t+\mu+\sigma\varepsilon_{t+1},
\qquad
\varepsilon_{t+1}\sim \mathcal N(0,1),
\]
where \(\mu\in\mathbb R\) is the drift and \(\sigma>0\) is the volatility. Let
\[
\mathcal F_t=\sigma(X_0,\ldots,X_t)
\]
be the information available at time \(t\).

\section{Deviation Signal and Predictive Funnel}

Fix a buying time \(\tau_b\). In the fixed-parameter baseline model, define the predictive center and cumulative variance by
\[
C(t)=X_{\tau_b}+\mu(t-\tau_b),
\qquad
V(t)=\sigma^2(t-\tau_b).
\]
For a threshold parameter \(k>0\), define the upper and lower funnel boundaries by
\[
U(t)=C(t)+k\sqrt{V(t)},
\qquad
L(t)=C(t)-k\sqrt{V(t)}.
\]
Equivalently, define the normalized deviation statistic
\[
Z_t=
\frac{X_t-X_{\tau_b}-\mu(t-\tau_b)}{\sigma\sqrt{t-\tau_b}},
\qquad t>\tau_b.
\]
Then \(X_t>U(t)\) is equivalent to \(Z_t>k\).

\begin{remark}[Fixed-time versus sequential interpretation]
For a fixed time \(t\), under the Gaussian model, \(Z_t\sim\mathcal N(0,1)\). However, the trading rule monitors \(Z_t\) repeatedly over time. Therefore the relevant object is not only the fixed-time tail probability \(P(Z_t>k)\), but the boundary-crossing probability
\[
P\left(\sup_{\tau_b<t\leq T}Z_t>k\right).
\]
Thus \(k\) should be interpreted as a funnel-width or boundary-crossing tuning parameter, not simply as a fixed-time significance level.
\end{remark}

\section{Momentum Signal}

For a window length \(\delta\geq 1\), define the finite-difference momentum signal
\[
M_t=X_t-X_{t-\delta}.
\]
When \(M_t>0\), the recent log-price movement is upward over the last \(\delta\) steps. When \(M_t\leq 0\), recent upward momentum has stalled or reversed. Under the baseline Gaussian model,
\[
M_t\sim\mathcal N(\mu\delta,\sigma^2\delta).
\]
Hence \(M_t\) is a noisy local trend signal.

\section{Transaction Costs and Profit Threshold}

Suppose that buying incurs proportional cost \(c_{\mathrm{buy}}\geq 0\) and selling incurs proportional cost \(c_{\mathrm{sell}}\in[0,1)\). Buying one unit at price \(S_{\tau_b}\) effectively costs
\[
(1+c_{\mathrm{buy}})S_{\tau_b},
\]
and selling at price \(S_t\) yields
\[
(1-c_{\mathrm{sell}})S_t.
\]
To avoid a realized loss after transaction costs, it is necessary and sufficient that
\[
(1-c_{\mathrm{sell}})S_t>(1+c_{\mathrm{buy}})S_{\tau_b}.
\]
Taking logarithms gives the profit condition
\[
X_t-X_{\tau_b}>h,
\]
where
\[
h=
\log\left(\frac{1+c_{\mathrm{buy}}}{1-c_{\mathrm{sell}}}\right).
\]
When transaction costs are ignored, \(h=0\).

\section{Stopping Rule}

\begin{definition}[Adaptive funnel stopping rule, fixed-parameter version]
Given a buying time \(\tau_b\), define the selling time by
\[
\tau_s=
\inf\left\{
 t>\tau_b:
 X_t>U(t),\quad M_t\leq 0,\quad X_t-X_{\tau_b}>h
\right\}.
\]
Equivalently,
\[
\tau_s=
\inf\left\{
 t>\tau_b:
 Z_t>k,\quad M_t\leq 0,\quad X_t-X_{\tau_b}>h
\right\}.
\]
\end{definition}

The rule says: sell only after the price crosses the upper funnel boundary, local momentum has weakened, and the trade is profitable after costs.

\section{Trading Cycle and Re-Entry Rule}

The stopping rule above describes when an existing long position is exited. To turn the rule into a repeated trading strategy, we also specify how the strategy re-enters after a completed sale.

Suppose the strategy sells at time \(\tau_s\) and then holds cash. During the cash period, take \(\tau_s\) as the new reference time and define
\[
Z_t^{\mathrm{down}}
=
\frac{X_t-X_{\tau_s}-\mu(t-\tau_s)}
{\sigma\sqrt{t-\tau_s}},
\qquad t>\tau_s.
\]
The next buying time is defined by the symmetric lower-funnel recovery rule
\[
\tau_b^{\mathrm{next}}
=
\inf\left\{
t>\tau_s:
Z_t^{\mathrm{down}}<-k,\quad M_t\geq 0
\right\}.
\]

Thus the full algorithm alternates between two regimes:
\begin{enumerate}
    \item while holding the asset, sell only when the upper-funnel, weakened-momentum, and profit-threshold conditions are all satisfied;
    \item while in cash, buy after a lower-funnel deviation once local momentum has recovered.
\end{enumerate}

The no-realized-loss theorem below applies to completed sell exits. The re-entry rule determines when capital is redeployed, but it does not by itself create a realized gain or loss.

\section{No-Realized-Loss Property}

\begin{theorem}[No realized loss on completed algorithmic exits]
Suppose a trade is entered at time \(\tau_b\) and exited at a finite time \(\tau_s<\infty\) according to the stopping rule above. Then the net realized wealth multiplier satisfies
\[
\frac{(1-c_{\mathrm{sell}})S_{\tau_s}}
{(1+c_{\mathrm{buy}})S_{\tau_b}}>1.
\]
Thus every completed algorithmic exit has positive realized return after transaction costs.
\end{theorem}

\begin{proof}
By the stopping rule,
\[
X_{\tau_s}-X_{\tau_b}>h
=
\log\left(\frac{1+c_{\mathrm{buy}}}{1-c_{\mathrm{sell}}}\right).
\]
Therefore
\[
\frac{S_{\tau_s}}{S_{\tau_b}}
=
\exp(X_{\tau_s}-X_{\tau_b})
>
\frac{1+c_{\mathrm{buy}}}{1-c_{\mathrm{sell}}}.
\]
Multiplying both sides by \((1-c_{\mathrm{sell}})/(1+c_{\mathrm{buy}})\) gives
\[
\frac{(1-c_{\mathrm{sell}})S_{\tau_s}}
{(1+c_{\mathrm{buy}})S_{\tau_b}}>1.
\]
\end{proof}

\begin{remark}
This theorem applies only to completed algorithmic exits satisfying the stopping rule. If a finite horizon \(T\) is imposed and the position is forcibly liquidated at \(T\), the theorem need not apply. Forced exits can realize losses.
\end{remark}

\section{Wealth Processes}

It is important to distinguish realized wealth from mark-to-market wealth. Let \(W_t^r\) denote realized wealth, updated only after completed exits. While holding a position purchased at \(\tau_b\), define mark-to-market wealth by
\[
W_t^m
=
W_{\tau_b}^r\frac{S_t}{S_{\tau_b}}
\]
in the frictionless case. With transaction costs interpreted as liquidation costs, one may instead use
\[
W_t^m
=
W_{\tau_b}^r
\frac{(1-c_{\mathrm{sell}})S_t}{(1+c_{\mathrm{buy}})S_{\tau_b}}.
\]
The no-realized-loss theorem concerns \(W_t^r\), not \(W_t^m\). While holding, \(W_t^m\) may decrease.

A natural benchmark is buy-and-hold wealth,
\[
W_t^{\mathrm{bh}}=W_0\frac{S_t}{S_0},
\]
which should be compared against both realized and mark-to-market strategy wealth.

\section{Momentum Decision Consistency}

The finite-difference signal \(M_t\) is a noisy local trend statistic. A simple consistency result can be obtained by working with returns. Suppose that over a local window,
\[
r_i=X_i-X_{i-1}=\beta+\epsilon_i,
\]
where \(\epsilon_i\) are independent with mean zero and finite variance. Define
\[
\hat\beta_n=\frac{1}{n}\sum_{i=1}^n r_i.
\]

\begin{theorem}[Consistency of local mean-return decision]
Under the assumptions above,
\[
\hat\beta_n\xrightarrow{p}\beta.
\]
Consequently, for any threshold \(c\neq \beta\),
\[
\mathbf 1\{\hat\beta_n>c\}
\xrightarrow{p}
\mathbf 1\{\beta>c\}.
\]
\end{theorem}

\begin{proof}
The convergence \(\hat\beta_n\xrightarrow{p}\beta\) follows from the law of large numbers. Since \(c\neq\beta\), the indicator map is continuous at \(\beta\), and the conclusion follows from the continuous mapping theorem.
\end{proof}

This theorem justifies a smoothed momentum rule based on average recent returns. The simpler finite-difference signal \(M_t=X_t-X_{t-\delta}\) corresponds to the cumulative return over the last \(\delta\) steps and is noisier, but follows the same intuition.

\section{Phase 2 Evaluation Metrics}

The Phase 2 simulation evaluates the strategy by separating accounting quantities that can otherwise be confused. The baseline simulation tracks three wealth curves:
\begin{enumerate}
    \item realized wealth \(W_t^r\), which updates only after completed exits;
    \item mark-to-market wealth \(W_t^m\), which shows hidden drawdowns while holding;
    \item buy-and-hold wealth \(W_t^{\mathrm{bh}}\), which measures opportunity cost.
\end{enumerate}

From these curves, the main diagnostic quantities are:
\[
R_T^r=\frac{W_T^r}{W_0}-1,\qquad
R_T^m=\frac{W_T^m}{W_0}-1,\qquad
R_T^{\mathrm{bh}}=\frac{W_T^{\mathrm{bh}}}{W_0}-1.
\]
The benchmark gap is measured by
\[
G_T=\frac{W_T^m}{W_T^{\mathrm{bh}}}-1,
\]
so \(G_T>0\) means the strategy's liquidating mark-to-market value exceeds buy-and-hold at the horizon, while \(G_T<0\) means the strategy underperforms buy-and-hold.

To measure hidden risk, define the mark-to-market drawdown process
\[
D_t^m=\frac{W_t^m}{\max_{0\leq s\leq t}W_s^m}-1,
\]
and its maximum drawdown over the horizon
\[
D_{\min}^m=\min_{0\leq t\leq T}D_t^m.
\]
The same calculation can be applied to \(W_t^{\mathrm{bh}}\) to compare the strategy's drawdown against the passive benchmark.

Finally, let \(H_t\in\{0,1\}\) denote whether the strategy is holding the asset. The exposure ratio
\[
E_T=\frac{1}{T+1}\sum_{t=0}^T H_t
\]
measures the fraction of time capital is exposed to the asset. Together, \(R_T^r\), \(R_T^m\), \(R_T^{\mathrm{bh}}\), \(G_T\), \(D_{\min}^m\), the number of completed trades, and \(E_T\) form the Phase 2 evaluation dashboard.

The simulation should use a Gaussian log-return model when validating the theory. If clipped shocks are used for visualization, they should be labelled as a visualization choice rather than as the exact model assumed in the paper.

\section{Discussion and Limitations}

The strategy enforces a no-realized-loss condition on completed voluntary exits. However, this does not imply that the strategy is risk-free. The main limitations are:
\begin{itemize}
    \item \textbf{Mark-to-market drawdown:} the current value of the position may fall while holding;
    \item \textbf{Indefinite holding:} if the stopping rule never triggers, the position may remain open indefinitely;
    \item \textbf{Forced-exit risk:} finite-horizon liquidation can realize losses;
    \item \textbf{Opportunity cost:} the strategy may underperform buy-and-hold;
    \item \textbf{Model misspecification:} real prices have changing volatility, jumps, and fat tails;
    \item \textbf{Parameter estimation error:} the drift \(\mu\) is difficult to estimate accurately.
\end{itemize}

Future phases should replace fixed \(\mu\) and \(\sigma\) with estimated or adaptive quantities, such as Kalman-filtered drift and GARCH volatility, yielding an adaptive funnel.

\section{Conclusion}

We formulated a funnel-based stopping strategy for trading under drift and volatility. The key correction is the explicit profit threshold \(h\), which gives a rigorous no-realized-loss theorem for completed algorithmic exits. The strategy remains risky in mark-to-market terms, and its practical value must be judged by simulation, buy-and-hold comparison, drawdown, and holding-time behavior.

\end{document}
